% Same Exposure, Opposite Signs -- SSRN working paper.
% Build: ./build.sh   (figures/tables from figs.py; Tectonic in conda env "tectonic")
% Numbers mirror ../DRAFT.md and the frozen outputs reconciled in ../NUMBERS_AUDIT.md.
\documentclass[11pt]{article}
\usepackage[letterpaper,margin=1.05in]{geometry}
\usepackage[T1]{fontenc}
\usepackage{newpxtext,newpxmath}
\usepackage{microtype}
\usepackage{setspace}\setstretch{1.12}
\usepackage{graphicx}
\usepackage{booktabs,tabularx,threeparttable}
\usepackage{float}
\usepackage[section]{placeins}
\usepackage{enumitem}
\usepackage[font=small,labelfont=bf,labelsep=period]{caption}
\usepackage{xcolor}
\usepackage[round,authoryear]{natbib}
\usepackage[colorlinks=true,linkcolor=black,citecolor=blue!50!black,urlcolor=blue!50!black]{hyperref}
\usepackage{xurl}
\newcommand{\doi}[1]{\href{https://doi.org/#1}{doi:#1}}
\usepackage{titlesec}
\titleformat{\section}{\large\bfseries}{\thesection.}{0.6em}{}
\titleformat{\subsection}{\normalsize\bfseries}{\thesubsection}{0.6em}{}
\setlist{itemsep=2pt,topsep=4pt}
\graphicspath{{figs/}}
\renewcommand{\topfraction}{0.85}\renewcommand{\bottomfraction}{0.6}
\renewcommand{\textfraction}{0.12}\renewcommand{\floatpagefraction}{0.75}
\newcommand{\confirm}[1]{\textcolor{red}{[#1]}}
\makeatletter\newcommand{\tabinput}[1]{\@@input #1 }\makeatother

\title{\textbf{Same Exposure, Opposite Signs:}\\[2pt]
\Large Interest-Rate Risk Reporting Conventions in SEC Form N-PORT}
\author{Rishabh Gyanendra Upadhyay\thanks{Inception42. E-mail: \href{mailto:uhrishabh@gmail.com}{uhrishabh@gmail.com}.
Code:\\ \href{https://github.com/HRishabh95/nport-dv100-sign-conventions}{github.com/HRishabh95/nport-dv100-sign-conventions}.
The views expressed are the author's own.}\\[2pt] \normalsize Inception42}
\date{October 2026}

\begin{document}
\maketitle
\thispagestyle{empty}

\begin{abstract}
\noindent The iShares Core U.S.\ Aggregate Bond ETF and the Vanguard Total Bond Market Index Fund hold
similar US investment-grade bonds, yet in their SEC filings one reports its interest-rate sensitivity
as positive and the other as negative. For a hypothetical \$100 million portfolio split equally
between them, adding the reported figures suggests almost no net sensitivity ($+\$0.09$ million);
harmonising the reporting conventions implies an estimated value change of about $-\$5.7$ million for
a parallel 100 basis point rise in rates. The cause is Form N-PORT itself, which asks for DV100, ``the
change in value of the portfolio resulting from a 100 basis point change in interest rates'', without
saying whether rates rise or fall. Using every public N-PORT filing from 2019Q4 to 2026Q3, we show
that 48\% of 53,603 plain long-only bond fund-quarters report a positive DV100 and 52\% a negative
one, in every quarter. The choice is made by the fund complex: 99.0\% of registrant-quarters with at
least three plain funds are unanimous, and a fund's sign is predicted by its siblings with 99.8\%
accuracy at 80.2\% coverage. A separate re-implementation reproduces the census. Read as filed,
49.2\% of cross-complex fund pairs show apparently offsetting rate exposures, against 4.0\% once signs
are harmonised. Secondary return analyses point the same way, while a pre-registered test of the
per-currency detail that a pending SEC proposal would remove is inconclusive. We provide code that
builds fund-complex convention lookup tables, and propose a one-sentence clarification to the form.

\medskip\noindent\textbf{Keywords:} Form N-PORT, DV100, interest-rate risk, regulatory reporting,
data comparability, mutual funds, ETFs.\\
\textbf{JEL classification:} G23, G28, G11.

\end{abstract}

\newpage
\setcounter{page}{1}

%=====================================================================
\section{Introduction}

Suppose a risk team holds two of the largest bond index funds in the United States, each tracking the
US investment-grade aggregate market, and wants to know the portfolio's exposure to a rise in rates.
Each fund reports the answer to the SEC every quarter: Item B.3 of Form N-PORT asks for DV100, the
change in portfolio value for a 100 basis point move in rates, by currency and maturity bucket. In
the filings for late 2025, the iShares Core U.S.\ Aggregate Bond ETF reports $+\$7.58$ billion and the
Vanguard Total Bond Market Index Fund reports $-\$21.02$ billion. Both hold investment-grade US bonds
with a duration near six years. One reports the gain from a fall in rates; the other reports the
loss from a rise. Added together at face value, an equal-weighted position in the two funds appears
almost insensitive to rates. Harmonised to one convention, its DV100 implies an estimated value change
of about $-5.7\%$ for a parallel 100 basis point rise (Figure~\ref{fig:example}).

\begin{figure}[!htb]
  \centering
  \includegraphics{fig0_example.pdf}
  \caption{\textbf{Same kind of bonds, opposite reported signs.} A hypothetical \$100m portfolio split
  equally between the iShares Core U.S.\ Aggregate Bond ETF and the Vanguard Total Bond Market Index Fund.
  Each segment is one fund's contribution to portfolio DV100, from the net DV100 each fund reported in the
  2025Q4 N-PORT data set (report dates August and September 2025). As filed, the two reported figures nearly
  cancel; on one convention (R, value change for a rise) both are negative. DV100-based sensitivity estimates
  for a parallel +100bp rise, not realised losses.}
  \label{fig:example}
\end{figure}

The cause is simple. The form defines DV100 as ``the change in value of the portfolio resulting from a
100 basis point change in interest rates'' \citep{sec_nport} without saying whether rates rise or
fall, and fund complexes have settled on different answers. This paper measures how widespread the
difference is, shows that it is a stable choice made by fund complexes rather than by individual
funds, and quantifies how much it changes mechanical comparisons of rate exposure across funds.

The question is timely. In February 2026 the SEC proposed amendments to Form N-PORT that would keep
DV100, require it in US dollars, and aggregate it across currencies \citep{sec_proposal}. Commenters
discussed the currency aggregation and the USD requirement; none of the 15 public comment letters
raised the direction of the shock \citep{sec_comments}. The industry's main trade association
supported the USD requirement precisely because it would ``facilitate comparisons across funds''
\citep{ici_letter}. Our evidence suggests that comparability also needs the direction stated.

We make four contributions. First, a full census: across 28 quarters and 105,416 fund-quarters, we
classify each filer's sign convention using only the filings themselves and contemporaneous
holdings, never returns, under a protocol frozen before the data were examined. Second, independent
validation: the inferred conventions agree with each fund's own holdings on 99.7\% of exposure, a
separate re-implementation reproduces the census, and the SEC bulk data match the original XML
filings cell for cell. Third, consequences: we show how harmonisation changes cross-fund comparisons
and, as secondary evidence, how it changes the explanation of subsequent fund returns. Fourth, a data
product: code that builds fund-complex and adviser lookup tables, which let users establish a filer's convention for any
quarter, together with a one-sentence clarification of the form.

\paragraph{Related literature.} This paper relates to three literatures. The first studies the
quality and comparability of mandatory regulatory data. Research on XBRL filings documents frequent
errors and inconsistencies in machine-readable disclosures \citep{debreceny2010,du2013}, and the
classic comparison of mutual fund databases shows that measurement choices in fund data can change
research conclusions \citep{elton2001}. Closest to our setting, \citet{flood2015} show that loosely specified risk fields in Form PF allow funds with different risks to report similar numbers. We document a different problem: a definitional ambiguity
that makes two correct filings incomparable. The second literature measures the interest-rate and
credit risk that bond funds take, and the fragility it creates \citep{becker2015,choi2018,goldstein2017,falato2021,ma2022};
this work relies on holdings-based or returns-based exposure measures \citep{barth2024}, and staff work that does use the N-PORT risk metrics adopts one reading of the sign: \citet{iorio2024} define DV01 as the value change for a one basis point \emph{decrease} in rates. Other recent work uses N-PORT holdings and derivatives \citep{kaniel2025}. Our results caution
against using the reported N-PORT risk metrics across complexes without harmonisation. The third
studies how funds use derivatives \citep{koski1999}; the funds whose rate exposure cannot be read
from holdings are exactly those where the reported metrics matter most. Methodologically, we follow
the recommendation to pre-register tests and report all outcomes \citep{harvey2017}, and we use
cluster bootstraps \citep{cameron2008} and an encompassing-style calibration test
\citep{chong1986,harvey1998}.

%=====================================================================
\section{Institutional background}\label{sec:background}

Form N-PORT, adopted in 2016 as part of the SEC's investment company reporting modernization
\citep{sec_2016}, requires registered funds other than money market funds to report portfolio
holdings and risk metrics monthly; the report for the third month of each fiscal quarter becomes
public. Funds that hold debt above a threshold (25\% of net assets) report portfolio-level risk
metrics in Item B.3: DV01 and DV100 for interest rates and CS01 for credit spreads, each by currency
and across five maturity buckets (3 months, 1, 5, 10 and 30 years), in the spirit of key-rate durations \citep{ho1992,tuckman2022}.

The instruction for Item B.3.b reads: ``provide the change in value of the portfolio resulting from a
100 basis point change in interest rates'' \citep{sec_nport}. Several other monetary items on the form
say ``Report in U.S.\ dollars''; Item B.3 says neither that nor which direction the shock goes, and
the adopting release adds neither. Two conventions are therefore consistent with the text:
\begin{itemize}
  \item \textbf{R (rise):} report the value change for a rise in rates, so a long bond fund reports a
  negative number.
  \item \textbf{F (fall or magnitude):} report a positive number for a long bond fund, either the value
  change for a fall in rates or an unsigned sensitivity.
\end{itemize}

The 2026 proposal would require risk metrics in US dollars, report DV100 aggregated across currencies,
drop DV01, merge investment-grade and non-investment-grade spread risk, and raise the reporting
threshold to 50\% of net assets in debt \citep{sec_proposal}. One comment letter argued that full
currency aggregation ``obscures localized, currency-specific convexity risks'' \citep{sharma_letter}.
No letter addressed the direction of the shock.

%=====================================================================
\section{Data}\label{sec:data}

We use the SEC's quarterly Form N-PORT public data sets from 2019Q4 to 2026Q3 \citep{sec_nportdata}, which contain reports for period-ends from September 2019 to July 2026 (quarters in this paper refer to these SEC data sets),
keeping the latest public report per fund (series) per data set. Table~\ref{tab:sample} summarises
the sample. Investment advisers come from Form N-CEN \citep{sec_ncen}. For the return analysis in
Section~\ref{sec:returns} we add monthly fund returns from N-PORT, official government yield curves
for six currencies, ECB reference exchange rates and factor ETF prices \citep{rates}.

\begin{table}[!htb]
\centering\small
\caption{Sample}\label{tab:sample}
\begin{tabular}{lr}
\toprule
N-PORT data sets (quarters) & 28 (2019Q4 to 2026Q3) \\
Fund-quarters reporting Item B.3 interest-rate fields & 105,416 \\
Distinct funds (series) & 6,073 \\
Distinct registrants & 1,531 \\
Distinct primary investment advisers (numeric CRD) & 509 \\
Plain long-only bond fund-quarters (definition in Section~\ref{sec:method}) & 53,603 \\
\quad of which used as convention evidence & 51,201 \\
Fund-quarters with a resolved convention & 86,276 (81.8\%) \\
Multi-currency offset fund-quarters (Section~\ref{sec:prevalence}) & 2,566 \\
\bottomrule
\end{tabular}
\end{table}

%=====================================================================
\section{Identifying the convention without returns}\label{sec:method}

The sign of a fund's DV100 is only informative if we know the true direction of its exposure. For
most funds we do not; for one large group we do. A \emph{plain} fund is a long-only bond fund: debt
is at least 50\% of net assets, rate derivatives are below 5\% of debt, short debt is below 1\%, and
an estimate of DV100 built from the fund's own holdings (each bond's duration from its coupon,
maturity and yield) is negative and at least 0.5\% of net assets. A 100 basis point rise in rates
certainly lowers the value of such a fund, so a negative reported DV100 reveals convention R and a
positive one convention F. Plain fund-quarters whose reported magnitude is outside one-third to three
times the holdings estimate (4.5\%) are not used as evidence.

Figure~\ref{fig:scatter} shows the result for all 51,201 evidence fund-quarters. If filers shared a
convention, reported DV100 would cluster along a single line through the holdings estimates. Instead
there are two mirror-image clouds: one along ``reported = estimate'' (convention R) and one along
``reported = $-$estimate'' (convention F). The magnitudes agree closely with the holdings estimates
in both clouds; only the sign differs.

\begin{figure}[!htb]
  \centering
  \includegraphics{fig1_scatter.pdf}
  \caption{\textbf{Two conventions for the same exposure.} Reported DV100 against the DV100 estimated
  from each fund's own holdings, both as \% of net assets, for 51,201 plain long-only bond
  fund-quarters, 2019Q4 to 2026Q3 (values clipped at $\pm$20\%). Every fund here loses value when rates
  rise (horizontal axis negative). Blue: reported negative (convention R); orange: reported positive
  (convention F). Dashed lines: reported = estimate and reported = $-$estimate.}
  \label{fig:scatter}
\end{figure}

Table~\ref{tab:example} makes this concrete with five index funds that track the same broad US
investment-grade bond market. All five have holdings-implied DV100 of about $-6\%$ of net assets.
Three report it as positive, two as negative.

\begin{table}[!htb]
\centering\footnotesize\setlength{\tabcolsep}{4pt}
\begin{threeparttable}
\caption{Worked example: aggregate bond index funds, 2025Q4 data set}\label{tab:example}
\begin{tabular}{lcrrrrc}
\toprule
Fund & Report date & NAV & \multicolumn{2}{c}{DV100 (\$bn)} & Reported & Conv. \\
\cmidrule(lr){4-5}
 & & (\$bn) & Reported & Holdings est. & (\% NAV) & \\
\midrule
\tabinput{tables/example.tex}
\bottomrule
\end{tabular}
\begin{tablenotes}\footnotesize
\item Reported: net Item B.3.b DV100 summed over currencies and buckets, as filed. Holdings estimate:
value change for a +100bp rise from the fund's own holdings. Illustration: a \$100m portfolio split
equally between the iShares and Vanguard funds has a face-value DV100 of $+\$0.09$m (apparently rate-neutral); with signs harmonised to convention R it is $-\$5.72$m.
\end{tablenotes}
\end{threeparttable}
\end{table}

Funds that are not plain (multi-asset, derivatives-heavy or multi-currency funds) cannot reveal their
own convention, because their true exposure may be of either sign. They inherit the convention of
their plain siblings in the same registrant, or failing that the same investment adviser, in the same
quarter, provided at least two siblings exist and at least 90\% of them agree. Records with no plain
sibling, with DV01 and DV100 signs that disagree within the same filing, or with unexplained unit
problems are left unresolved rather than assigned a sign. Local-currency reporting is detected where
the ratio of a non-USD leg to its holdings estimate matches the exchange rate (493 legs, mostly JPY)
and converted. The full rule set was frozen before the audit was run; one labelled deviation (D1),
made afterwards, keeps as convention R the 12,703 records whose DV100 exceeds 100 times DV01, reflecting callable municipal bonds' negative convexity.

%=====================================================================
\section{How common is each convention?}\label{sec:prevalence}

The split is not a transient or a small-sample artefact. Figure~\ref{fig:quarters} shows the share
of plain fund-quarters reporting a positive DV100 in each of the 28 data sets: between 43\% and 50\%
every quarter, with no trend. Over the whole sample, 25,735 of 53,603 plain fund-quarters (48\%)
report positive.

\begin{figure}[!htb]
  \centering
  \includegraphics{fig2_quarters.pdf}
  \caption{\textbf{About half of plain bond funds report a positive DV100, every quarter.} Share of plain
  long-only bond fund-quarters reporting positive (orange) and negative (blue) net DV100, by N-PORT
  data set. Dotted line: 50\%.}
  \label{fig:quarters}
\end{figure}

The convention follows the fund complex. Figure~\ref{fig:complexes} plots, for the largest registrants of each convention, the share of their plain fund-quarters reporting positive. Almost
every registrant sits at 0\% or 100\%. iShares, Fidelity, Northern, DoubleLine and PIMCO report
positive; DFA, SPDR, Invesco, Franklin and Eaton Vance report negative. Across all registrant-quarters
with at least three plain funds, 99.0\% are unanimous, and a plain fund's convention is predicted by
the majority of its plain siblings with 99.8\% accuracy, covering 41,074 of the 51,201 evidence
fund-quarters (80.2\%).

\begin{figure}[!htb]
  \centering
  \includegraphics{fig3_complexes.pdf}
  \caption{\textbf{Within a fund complex, almost everyone agrees.} Share of each registrant's plain
  fund-quarters reporting a positive (orange) or negative (blue) DV100, for the eight registrants with the
  most plain fund-quarters in each convention; plain fund-quarters at right. Allspring changed convention in
  2023Q1. Table~\ref{tab:complexes} and the lookup tables give the full evidence.}
  \label{fig:complexes}
\end{figure}

Table~\ref{tab:complexes} lists the largest registrants with their dominant convention and how
consistently they apply it.

\begin{figure}[!tb]
  \centering
  \includegraphics{fig4_status.pdf}
  \caption{\textbf{How much of N-PORT DV100 is comparable as filed.} Status of each fund-quarter after
  the audit (with deviation D1), by count and by net assets, for all 105,416 fund-quarters and for the
  2,566 multi-currency offset fund-quarters (cross-currency offset at least 0.20, non-USD share at least
  0.10). Unresolved records are never assigned a sign.}
  \label{fig:status}
\end{figure}

\begin{table}[!htb]
\centering\small
\begin{threeparttable}
\caption{Conventions of the largest registrants}\label{tab:complexes}
\begin{tabular}{lrcrrr}
\toprule
Registrant & Plain fund- & Dominant & Share & Quarters & Changes \\
 & quarters & conv. & (\%) & resolved & \\
\midrule
\tabinput{tables/complexes.tex}
\bottomrule
\end{tabular}
\begin{tablenotes}\footnotesize
\item Twelve registrants with the most plain fund-quarters. Share: fraction of plain fund-quarters with
the dominant convention. Quarters resolved: registrant-quarters classified R or F in the lookup table
(at least two plain funds, at least 90\% agreement). Changes: switches between resolved quarters.
\end{tablenotes}
\end{threeparttable}
\end{table}

Conventions are stable but not permanent. Twenty-six registrants changed convention between resolved
quarters (Table~\ref{tab:changes}), and when they change, the whole complex changes at once. Allspring
reported positive through 2022Q4 and negative from 2023Q1; several Nuveen trusts switched to positive
in 2022Q1 and back in 2023Q1. A user who maps each complex to a single convention will therefore be
wrong for some quarters; the convention has to be read by quarter.

\begin{table}[!htb]
\centering\small
\caption{Registrants that changed convention (first 10 of 31 changes, by quarter)}\label{tab:changes}
\begin{tabular}{lccr}
\toprule
Registrant & Data set & Change & Plain funds \\
\midrule
\tabinput{tables/changes.tex}
\bottomrule
\end{tabular}
\end{table}

What do the positive numbers mean? Fewer than 1\% of filings explain their convention in an
explanatory note. The two complexes that do support the reading above: MassMutual describes an average
of up and down shifts, i.e.\ a magnitude, and Ellington describes ``a 100 bp decrease''. Among F
filers the median ratio of DV100 to 100 times DV01 is exactly 1.000, so most appear to scale DV01 up
rather than revalue the portfolio at 100 basis points.

\paragraph{How much of the data can be compared.} Figure~\ref{fig:status} shows the status of every
fund-quarter. Across all funds, 42\% of fund-quarters are comparable as filed (convention R, US
dollars), 40\% become comparable after a supported sign flip or unit conversion, and 18\% (13\% of
net assets) remain unresolved. For multi-currency funds whose currency legs offset one another, the
funds most affected by the proposal, only 20\% of net assets is comparable as filed; 65\% becomes comparable after a sign flip or unit conversion, and 14\% remains unresolved.


%=====================================================================
\section{Validation}\label{sec:validation}

Because the inference for non-plain funds borrows a sibling's convention, we test it with evidence the
inference did not use. Table~\ref{tab:validation} summarises four checks.

\begin{table}[!htb]
\centering\small
\begin{threeparttable}
\caption{Validation of the convention classification}\label{tab:validation}
\begin{tabularx}{\linewidth}{>{\raggedright\arraybackslash}p{4.2cm}X}
\toprule
Check & Result \\
\midrule
Leave-one-out on plain funds & Sibling majority predicts 99.8\% correctly; covers 41,074 of 51,201
evidence fund-quarters (80.2\%). \\
Inferred conventions vs.\ each fund's own holdings (frozen protocol) & On 45,467 currency legs in 29,552
of 35,645 inferred fund-quarters, the harmonised sign agrees with the holdings estimate on 99.7\% of
exposure (98.3\% of legs). 6,093 inferred fund-quarters have no material leg and are unchecked;
derivative exposure the holdings model cannot value is not tested. \\
Separate re-implementation & An independent model run, with its own stricter definitions on the same
SEC bulk data and no access to the audit code, finds 45.5\% positive (41 to 48\% by quarter), 99.4\%
unanimity, leave-one-out accuracy 99.9\% at 84.0\% coverage, and the same dominant sign for each of
the 15 largest registrants. \\
Bulk data vs.\ original filings & For 19 filings drawn by rule, all 950 DV01/DV100 cells in the SEC
bulk data match the original XML exactly. \\
\bottomrule
\end{tabularx}
\begin{tablenotes}\footnotesize
\item Five families fail the holdings check at family level (three multi-adviser umbrella trusts, Stone
Ridge and Stone Harbor). Reclassifying their 68 records, and 15 records inferred through an N-CEN
placeholder adviser identifier, as unresolved leaves every reported figure unchanged.
\end{tablenotes}
\end{threeparttable}
\end{table}

%=====================================================================
\section{Consequences for cross-fund comparison}\label{sec:consequences}

The most direct consequence is for anyone who sums, nets or ranks DV100 across funds from different
complexes. Figure~\ref{fig:pairs}(a) takes every pair of resolved fund-quarters from different
registrants in the same quarter. As filed, 49.2\% of the 97 million pairs have net DV100s of
opposite sign: read as filed, their rate exposures appear to offset in direction, although with different magnitudes they need not be close to a hedge. After harmonising signs, 4.0\% do; these are genuine directional offsets, such as a fund with net short duration paired with a conventional bond fund. Under deviation D1,
which resolves more records, the figures are 50.2\% and 3.4\% of 133 million pairs.

\begin{figure}[!htb]
  \centering
  \includegraphics{fig5_pairs.pdf}
  \caption{\textbf{Harmonisation changes cross-fund comparisons.} (a) Share of cross-complex pairs of
  resolved fund-quarters whose net DV100s have opposite signs, as filed and with signs harmonised
  (frozen audit rule; unweighted; units not converted). (b) Share of 23,566 cross-family pairs of
  multi-currency funds whose diversification class (offsetting: predicted rate-P\&L correlation
  $<-0.3$; similar: $>+0.3$; otherwise neutral; 60-month yield covariance) changes between raw and
  harmonised data, and between the proposal's currency aggregate and full per-currency detail.
  Model-based classifications; neither change is validated by realised returns.}
  \label{fig:pairs}
\end{figure}

For multi-currency funds, a finer question is whether two funds' rate exposures are offsetting,
neutral or similar. Classifying 23,566 cross-family pairs by their predicted rate-P\&L correlation,
harmonising signs and units changes the class of 48.9\% of pairs; moving from the proposal's
currency aggregate to full per-currency detail changes 18.9\% (Figure~\ref{fig:pairs}(b)). These are
model-based classifications. Where the aggregate and the detail predict opposite-signed
correlations, realised residual correlations sided with the detail 186 times and with the aggregate
180 times, so the pair-level classification is not validated by returns. We have no evidence about
actual investor decisions or losses.

%=====================================================================
\section{Evidence from returns}\label{sec:returns}

Do the differences matter for explaining what funds subsequently earn? We answer this as secondary
evidence, with a design pre-registered before any return was read and amended twice, with dated
amendments, before outcomes were examined. We compare identically cleaned information sets: A
(holdings and derivatives only); B (A plus the proposal-equivalent aggregate DV100, used with each
holding's currency as the SEC describes); C (A plus harmonised per-currency DV100); and
$\mathrm{B}_{\mathrm{raw}}$ and $\mathrm{C}_{\mathrm{raw}}$ (B and C as filed). Each record is used
only from the month after its filing date. Rate P\&L is priced at realised monthly yield changes,
and the remaining return is explained by a fund-specific rolling 36-month factor model refitted for
each set. The outcome is the error of this \emph{conditional} prediction; it is not a forecast and not
a tradable signal. Confidence intervals are the more conservative of a calendar-month and a
fund-family cluster bootstrap \citep{cameron2008}. Development and confirmation samples are split by
fund family, and funds inspected earlier are excluded from confirmation.

\begin{table}[!htb]
\centering\small
\begin{threeparttable}
\caption{Return evidence (confirmation sample)}\label{tab:returns}
\begin{tabularx}{\linewidth}{>{\raggedright\arraybackslash}X>{\raggedright\arraybackslash}p{3.9cm}>{\raggedright\arraybackslash}p{3.9cm}}
\toprule
 & \textbf{Primary, pre-registered} (153 funds, 2,117 fund-months) & \textbf{Broad, secondary} (343
 funds, 10,988 fund-months) \\
\midrule
Frozen calibration test $\theta$ (passes at $\geq 0.5$) & \textbf{0.496 [0.135, 0.992]: inconclusive} &
0.72 [0.41, 1.02] \\
Per-currency detail: MSE reduction, C vs.\ B & 6.7\% [$-1.3$, 32.4] & 0.9\% [0.25, 1.5] \\
Signs and units harmonised: MSE reduction, B vs.\ $\mathrm{B}_{\mathrm{raw}}$ (descriptive) &
10.0\% [$-1.2$, 13.8]$^\dagger$ & 9.7\% [4.7, 13.7]$^\dagger$ \\
\bottomrule
\end{tabularx}
\begin{tablenotes}\footnotesize
\item 95\% intervals in brackets: the more conservative of month and family bootstraps. MSE: conditional
prediction error of monthly fund returns. Primary population: multi-currency offset funds; broad
population: funds with at least 10\% non-USD debt. $\theta$: slope of B's prediction error on the
difference between C's and B's rate P\&L.
\item[$\dagger$] Interval computed after the frozen run. Sign and unit effects are not separated.
\end{tablenotes}
\end{threeparttable}
\end{table}

Table~\ref{tab:returns} reports the results. The pre-registered primary test, in the population the
proposal most affects, is inconclusive: $\theta$ falls 0.004 short of its threshold, and one
managed-futures fund supplies 89\% of its signal (its holdings estimate swings by $\pm$20\% of net
assets a month, and its reported per-currency legs correct it). Without that fund, a post-hoc
calculation, per-currency detail improves prediction error by 0.8\% [$-2.2$, 2.2]. In the broader
secondary population, harmonising signs and units lowers prediction error by 9.7\%, and adding
per-currency detail by 0.9\% (Figure~\ref{fig:returns}); the currency-detail gain is absent in
2025 and 2026. The two are not like-for-like: a user can repair the sign convention from public
information for most records, whereas currency detail removed by aggregation cannot be recovered.

\begin{figure}[!htb]
  \centering
  \includegraphics{fig6_returns.pdf}
  \caption{\textbf{Secondary analysis: harmonisation vs.\ currency detail.} Reduction in the conditional
  prediction error (MSE) of monthly fund returns, broad population (10,988 fund-months, 343 funds with
  at least 10\% non-USD debt, confirmation sample), with 95\% intervals. $^\dagger$Interval computed
  after the frozen run; sign and unit effects are not separated. The pre-registered primary test was
  inconclusive (Table~\ref{tab:returns}).}
  \label{fig:returns}
\end{figure}

%=====================================================================
\section{Implications}\label{sec:implications}

\paragraph{For users of N-PORT data.} Risk teams, allocators and data vendors should not sum, net or
rank DV100 across fund complexes without establishing each complex's convention for that quarter. The
lookup tables built by the accompanying code resolve 76.9\% of fund-quarters, trying the registrant first and then the
adviser; the remainder should be treated as unknown. An F label describes how a complex reported, not
an instruction to flip every number: it combines fall-direction and magnitude reporting, which cannot
be distinguished from the numbers alone.

\paragraph{For the form.} One sentence would remove the ambiguity, for example: ``report the change in
value for a parallel 100 basis point increase in interest rates; a decrease in value is reported as a
negative number.'' The proposed USD requirement fixes units, not direction.

\paragraph{For the proposed currency aggregation.} This paper does not show that aggregation harms
risk assessment on average. The cost appears concentrated in a few derivatives-heavy funds, and that
conclusion rests on post-hoc evidence from one fund.

\paragraph{Limitations.} DV100 is treated as linear and parallel by bucket; flipping F records ignores
convexity asymmetry (about 1 to 5\% for plain bond funds), and magnitude-style figures are assumed to
have the economic sign of a rise. Positions are 1 to 9 months stale. Rates are mapped for six
currencies; the rest are priced on each fund's currency mix. Factor returns exclude ETF distributions.
The analysis covers public quarterly N-PORT reports only, says nothing about the proposed threshold
change, the IG/non-IG merge or the removal of DV01, and does not measure investor behaviour or losses.

%=====================================================================
\section{Conclusion}

A single number in a mandatory SEC filing means opposite things depending on who files it. About half
of plain bond funds report the interest-rate sensitivity of their portfolios as positive and half as
negative, and the choice is made consistently at the level of the fund complex. Taken at face value,
these figures make about half of all cross-complex fund pairs look like offsetting rate exposures.
The problem is easy to fix for future filings with one sentence in the form instructions, and for
past filings with the lookup tables that the accompanying code builds.

\section*{Data and code availability}
All inputs are public. Form N-PORT and Form N-CEN data sets are available from the SEC; yield curves and
exchange rates from the official sources listed in the references. The code for the audit, the lookup
tables, the outcome study, the separate re-implementation and this paper's figures and tables is
available at \url{https://github.com/HRishabh95/nport-dv100-sign-conventions}, and regenerates every
reported number from those public sources. Protocol hashes: audit 7c51d52c, outcome fe1c473e,
amendment~1 e6a248e8, amendment~2 8ae3f579, sign check 29cbe65b. Labelled deviations: D1
(Section~\ref{sec:method}); an exclusion-matching bug fixed after the development run (development
numbers unchanged); post-hoc intervals for the harmonisation comparison; post-hoc leave-one-fund-out
figures.

\section*{Use of generative AI}
Anthropic's Claude Opus 5.5 models were used to develop the analysis code; extract and analyse the
public SEC data; generate the figures; draft the manuscript text; and run, as separate model
instances, the independent re-implementation of the census and an overclaiming review of the
manuscript. The author selected the research question, approved the research protocols and their
amendments, reviewed the evidence and conclusions, directed revisions, and accepts responsibility for
the final manuscript.

\begin{thebibliography}{99}\small\setlength{\itemsep}{2pt}
\bibitem[{Barth et al.}(2024)]{barth2024} Barth, D., R.\,J. Kahn, P. Monin, and O. Sokolinskiy (2024). Reaching for Duration and Leverage in the Treasury Market. Finance and Economics Discussion Series 2024-039, Board of Governors of the Federal Reserve System. \doi{10.17016/FEDS.2024.039}
\bibitem[{Becker and Ivashina}(2015)]{becker2015} Becker, B., and V. Ivashina (2015). Reaching for Yield in the Bond Market. \emph{Journal of Finance} 70(5), 1863-1902. \doi{10.1111/jofi.12199}
\bibitem[{Cameron et al.}(2008)]{cameron2008} Cameron, A.\,C., J.\,B. Gelbach, and D.\,L. Miller (2008). Bootstrap-Based Improvements for Inference with Clustered Errors. \emph{Review of Economics and Statistics} 90(3), 414-427. \doi{10.1162/rest.90.3.414}
\bibitem[{Choi and Kronlund}(2018)]{choi2018} Choi, J., and M. Kronlund (2018). Reaching for Yield in Corporate Bond Mutual Funds. \emph{Review of Financial Studies} 31(5), 1930-1965. \doi{10.1093/rfs/hhx132}
\bibitem[{Chong and Hendry}(1986)]{chong1986} Chong, Y.\,Y., and D.\,F. Hendry (1986). Econometric Evaluation of Linear Macro-Economic Models. \emph{Review of Economic Studies} 53(4), 671-690. \doi{10.2307/2297611}
\bibitem[{Debreceny et al.}(2010)]{debreceny2010} Debreceny, R., S. Farewell, M. Piechocki, C. Felden, and A. Gr\"aning (2010). Does It Add Up? Early Evidence on the Data Quality of XBRL Filings to the SEC. \emph{Journal of Accounting and Public Policy} 29(3), 296-306. \doi{10.1016/j.jaccpubpol.2010.04.001}
\bibitem[{Du et al.}(2013)]{du2013} Du, H., M.\,A. Vasarhelyi, and X. Zheng (2013). XBRL Mandate: Thousands of Filing Errors and So What? \emph{Journal of Information Systems} 27(1), 61-78. \doi{10.2308/isys-50399}
\bibitem[{Elton et al.}(2001)]{elton2001} Elton, E.\,J., M.\,J. Gruber, and C.\,R. Blake (2001). A First Look at the Accuracy of the CRSP Mutual Fund Database and a Comparison of the CRSP and Morningstar Mutual Fund Databases. \emph{Journal of Finance} 56(6), 2415-2430. \doi{10.1111/0022-1082.00410}
\bibitem[{Falato et al.}(2021)]{falato2021} Falato, A., I. Goldstein, and A. Hortaçsu (2021). Financial Fragility in the COVID-19 Crisis: The Case of Investment Funds in Corporate Bond Markets. \emph{Journal of Monetary Economics} 123, 35-52. \doi{10.1016/j.jmoneco.2021.07.001}
\bibitem[{Flood et al.}(2015)]{flood2015} Flood, M.\,D., P. Monin, and L. Bandyopadhyay (2015). Gauging Form PF: Data Tolerances in Regulatory Reporting on Hedge Fund Risk Exposures. OFR Working Paper 15-13, Office of Financial Research. \url{https://www.financialresearch.gov/working-papers/2015/07/30/gauging-form-pf/}
\bibitem[{Goldstein et al.}(2017)]{goldstein2017} Goldstein, I., H. Jiang, and D.\,T. Ng (2017). Investor Flows and Fragility in Corporate Bond Funds. \emph{Journal of Financial Economics} 126(3), 592-613. \doi{10.1016/j.jfineco.2016.11.007}
\bibitem[{Harvey}(2017)]{harvey2017} Harvey, C.\,R. (2017). Presidential Address: The Scientific Outlook in Financial Economics. \emph{Journal of Finance} 72(4), 1399-1440. \doi{10.1111/jofi.12530}
\bibitem[{Harvey et al.}(1998)]{harvey1998} Harvey, D.\,I., S.\,J. Leybourne, and P. Newbold (1998). Tests for Forecast Encompassing. \emph{Journal of Business \& Economic Statistics} 16(2), 254-259. \doi{10.1080/07350015.1998.10524759}
\bibitem[{Ho}(1992)]{ho1992} Ho, T.\,S.\,Y. (1992). Key Rate Durations: Measures of Interest Rate Risks. \emph{Journal of Fixed Income} 2(2), 29-44. \doi{10.3905/jfi.1992.408049}
\bibitem[{Investment Company Institute}(2026)]{ici_letter} Investment Company Institute (2026).
Comment letter on File No.\ S7-2026-05, 24 April 2026.
\url{https://www.sec.gov/comments/S7-2026-05/s7202605-759947-2336774.pdf}
\bibitem[{Iorio et al.}(2024)]{iorio2024} Iorio, B., D. Li, and L. Petrasek (2024). Why Do Mutual Funds Invest in Treasury Futures? FEDS Notes, Board of Governors of the Federal Reserve System, 10 May 2024. \doi{10.17016/2380-7172.3488}
\bibitem[{Kaniel and Wang}(2025)]{kaniel2025} Kaniel, R., and P. Wang (2025). Unmasking Mutual Fund Derivative Use. \emph{Review of Financial Studies} 38(4), 1120-1166. \doi{10.1093/rfs/hhaf001}
\bibitem[{Koski and Pontiff}(1999)]{koski1999} Koski, J.\,L., and J. Pontiff (1999). How Are Derivatives Used? Evidence from the Mutual Fund Industry. \emph{Journal of Finance} 54(2), 791-816. \doi{10.1111/0022-1082.00126}
\bibitem[{Ma et al.}(2022)]{ma2022} Ma, Y., K. Xiao, and Y. Zeng (2022). Mutual Fund Liquidity Transformation and Reverse Flight to Liquidity. \emph{Review of Financial Studies} 35(10), 4674-4711. \doi{10.1093/rfs/hhac007}
\bibitem[{Sharma}(2026)]{sharma_letter} Sharma, R. (2026). Comment letter on File No.\ S7-2026-05,
15 March 2026. \url{https://www.sec.gov/comments/s7-2026-05/s7202605-724347-2266094.pdf}
\bibitem[{Tuckman and Serrat}(2022)]{tuckman2022} Tuckman, B., and A. Serrat (2022). \emph{Fixed Income Securities: Tools for Today's Markets}, 4th ed. Hoboken, NJ: Wiley. \doi{10.1002/9781119835622}
\bibitem[{SEC}(2016)]{sec_2016} U.S. Securities and Exchange Commission (2016). Investment Company
Reporting Modernization. Release No.\ 33-10231.
\url{https://www.sec.gov/files/rules/final/2016/33-10231.pdf}
\bibitem[{SEC}(2026a)]{sec_proposal} U.S. Securities and Exchange Commission (2026a). Form N-PORT
Reporting. Proposed rule, Release No.\ IC-35962, File No.\ S7-2026-05.
\url{https://www.sec.gov/files/rules/proposed/2026/ic-35962.pdf}
\bibitem[{SEC}(2026b)]{sec_comments} U.S. Securities and Exchange Commission (2026b). Public comments
on File No.\ S7-2026-05. \url{https://www.sec.gov/rules-regulations/public-comments/s7-2026-05}
\bibitem[{SEC}(n.d.a)]{sec_nport} U.S. Securities and Exchange Commission (n.d.a). Form N-PORT:
Monthly Portfolio Investments Report (form and instructions). \url{https://www.sec.gov/files/formn-port.pdf}
\bibitem[{SEC}(n.d.b)]{sec_nportdata} U.S. Securities and Exchange Commission (n.d.b). Form N-PORT
Data Sets, 2019Q4 to 2026Q3. \url{https://www.sec.gov/data-research/sec-markets-data/form-n-port-data-sets}
\bibitem[{SEC}(n.d.c)]{sec_ncen} U.S. Securities and Exchange Commission (n.d.c). Form N-CEN Data
Sets. \url{https://www.sec.gov/data-research/sec-markets-data/form-n-cen-data-sets}
\bibitem[{Rate sources}(n.d.)]{rates} Yield and exchange-rate sources: U.S. Treasury, Daily Treasury Par
Yield Curve Rates; European Central Bank Data Portal, euro area AAA spot curve (series YC) and euro
foreign exchange reference rates (series EXR); Bank of England, nominal government spot curve; Japan
Ministry of Finance, JGB interest rates; Bank of Canada, Valet API; Reserve Bank of Australia,
statistical tables F1 and F2.
\end{thebibliography}

\end{document}
